\chapter{Бесцветный Ходж-вакуум и калибровочный якорь}
\label{ch:v8-colorless-hodge-gauge-anchor-no-go-gate}

\section{Последний допустимый носитель}

После всех цветовых проверок Том VII оставил в активном Ходж-вакууме ровно
четыре ребра
\begin{equation}
 \mathcal E_-=\{
 \mathrm{L_L-Y_R},\ \mathrm{X_L-X_R},\
 \mathrm{X_L-e_R},\ \mathrm{Y_L-Y_R}
 \}.
\end{equation}
Каждое из них несёт представление $(\mathbf1,\mathbf1)_0$. Поэтому сумма
калибровочных индексов активного сектора тождественно равна
\begin{equation}
 (I_1,I_2,I_3)_{\mathcal E_-}=(0,0,0).
 \label{eq:v8-active-gauge-index-zero}
\end{equation}
При этом Ходж-потенциал на тех же четырёх рёбрах ненулевой и именно он
создаёт конденсат. Иными словами, активный скалярный сектор существует, но
для калибровочной кинетики он невидим.

\section{Почему полный след не исправляет нуль}

Все индексы полного одиннадцатирёберного носителя,
\begin{equation}
 (I_1,I_2,I_3)_{\rm full}=\left(13,2,\frac32\right),
\end{equation}
целиком находятся на семи положительных, не конденсирующихся рёбрах.
Следовательно, условная формула прежнего гейта
\begin{equation}
 f_0=\frac{6\pi^2}{q_Gg^2},\qquad
 \lambda_E=\frac{q_G}{6}g^2
\end{equation}
не может быть применена к активному Ходж-сектору: на нём $q_G=0$ для
каждого калибровочного фактора.

Можно вернуть в след семь заряженных невзаимодействующие степени свободы, но тогда калибровочная
нормировка берётся уже с другого носителя. Кроме того, три индекса
$13,2,3/2$ не являются одним универсальным числом. Без явно построенного
общего оператора произведения это не фиксирует относительный вес между
калибровочной кривизной и бесцветным Ходж-потенциалом.

\begin{theorem}[Калибровочная невидимость физического Ходж-вакуума]
Окончательный цветосохраняющий Ходж-вакуум Тома VII имеет нулевые индексы
относительно $U(1)_Y$, $SU(2)_w$ и $SU(3)_c$. Поэтому его собственный
калибровочный след не может определить $f_0$ или массовую метрику.
Ненулевой калибровочный след полного рёберного носителя целиком создаётся
заглушённым сектором и без общего оператора произведения не замыкает тот же
скалярный функционал.
\end{theorem}

\begin{nogocriterion}[Запрет формального переноса $q_G$]
Нельзя подставлять полный индекс одиннадцати рёбер или старое значение
$q_G=2$ в формулу квартики четырёхрёберного вакуума. Такой перенос смешивает
два разных следовых носителя и скрыто выбирает искомую относительную метрику.
\end{nogocriterion}

\begin{protocol}
\begin{enumerate}
 \item окончательная активная опора: \textbf{четыре калибровочного синглета};
 \item её индексы $(I_1,I_2,I_3)$: \textbf{$(0,0,0)$};
 \item полный индекс одиннадцати рёбер: \textbf{$(13,2,3/2)$};
 \item носитель полного индекса: \textbf{семь заглушённых рёбер};
 \item активный калибровочный якорь $f_0$: \textbf{отсутствует};
 \item единый калибровочно-ходжев след: \textbf{не выведен};
 \item ветвь общего калибровочно-пространственный следа в текущей архитектуре:
 \textbf{закрыта};
 \item оставшийся подготовительный примитив: \textbf{второй независимый
 семейный тензор}.
\end{enumerate}
\end{protocol}

Машинный сертификат сохранён в
\path{s2t_v8_colorless_hodge_gauge_anchor_no_go_gate_results.json}.